\chapter{Background and Related Work}
\label{chap:background-related-work}

This chapter connects several bodies of work needed to construct an observational causal-analysis testbed: causal-method evaluation, critical-care data and irregular time series, electronic phenotyping and programmatic labeling, bounded LLM-assisted design, and observational causal inference.  The emphasis is comparative: each supplies a component or boundary, but none alone defines the complete resource constructed here.  Project-specific implementation begins in Chapter~\ref{chap:problem-definition-study-design}.

\section{Causal-Estimator Evaluation: Control and Observational Fidelity}
\label{sec:background-related-work:causal-estimator-evaluation-spectrum}

Evaluating a causal estimator differs from evaluating a factual-outcome predictor because the individual counterfactual outcome is not observed.  Evaluation designs therefore differ in how they obtain an answer key and in how much of the observed system they preserve.  The resulting spectrum is not a ranking in which one setting dominates the others: control supports direct error measurement, whereas observational fidelity preserves difficulties that a specified mechanism may omit.

\subsection{Controlled answer-key designs}

Fully synthetic studies specify covariates, exposure assignment, potential outcomes, confounding, support, and effect heterogeneity, making estimator error measurable under evaluator-chosen mechanisms.  Semi-synthetic studies retain empirical components while constructing assignments or outcomes; an IHDP-style design illustrates this trade-off \cite{shalit2017ite}.  Experimentally anchored evaluation can instead induce an observational comparison from randomized data and compare it with the experimental reference \cite{gentzel2021experimental}.  Empirically anchored generation preserves additional distributional structure while still specifying the causal mechanisms and answer key.  Across these designs, control enables error measurement, but results remain conditional on construction, sampling, population, and harmonization choices.

\subsection{Observational testbeds from source-recorded data}

A source-recorded observational testbed retains routine measurements, availability and missingness patterns, care-process traces, covariates, and outcomes.  Cohort selection, proxy definitions and prevalence, temporal cutoffs, aggregation, and empirical support arise from the researcher-defined representation.  Controlled benchmarks provide measurable causal accuracy under specified mechanisms; CliniCause instead exposes estimator behavior under source-recorded empirical structure and explicit representation choices while the actual causal mechanism and truth remain unknown.  The modes are complementary rather than interchangeable.

\section{ICU EHR Datasets and Irregular Sampling}
\label{sec:background-related-work:icu-ehr-datasets-irregular-sampling}

\subsection{Irregular measurement and observation}

ICU records combine physiology, laboratories, treatments, observations, and administrative information without a common sampling clock.  Irregular spacing within variables, heterogeneous rates across variables, and patient-specific measurement density challenge complete fixed-dimensional sequence models \cite{sun_2026_review_irregular_medical_timeseries}.  Absence can reflect a decision not to measure, device or table availability, an analysis window, or preprocessing; masks and elapsed time can carry predictive information without identifying the process that generated it \cite{che2018grud}.  Identifiers, time zero, cohort entry, repeated stays, and summary windows likewise determine the analytical unit and temporal interpretation.

Regular resampling and imputation choose a grid and a rule for unobserved cells.  Sparse-event representations preserve observed values and times, while missingness-aware recurrence encodes masks and gaps.  Neither removes the need to distinguish value, observation, and care processes, so irregular sampling remains both a representation problem and an interpretation boundary.

\subsection{PhysioNet 2012 and MIMIC-III}

The PhysioNet/Computing in Cardiology Challenge 2012 provides a bounded early-ICU resource for mortality prediction, with admission descriptors and time-stamped vital and laboratory series \cite{silva2012physionet}.  The source documents the challenge, not the exact later CliniCause inclusion rules, missing-value handling, splits, or scores.

MIMIC-III is a deidentified, single-center relational critical-care database spanning bedside observations, laboratories, medications, procedures, coded information, notes, and outcomes \cite{johnson2016mimiciii}.  Investigators must define its study unit, cohort, variables, windows, and outcome.  Benchmark work demonstrates standardized task construction but is a particular extraction rather than the database or the local proxy task \cite{harutyunyan_2019_mimiciii_benchmark}.

PhysioNet's bounded challenge role and MIMIC-III's broader relational role test workflow portability across different scales, eras, identifiers, variables, and measurement systems.  Dataset-specific cohorts, rules, and graphs are retained; similarly named states need not be equivalent, and no populations or estimates are pooled.

\section{Irregular Time-Series Representation Models}
\label{sec:background-related-work:irregular-time-series-representation-models}

\subsection{Representation challenges and design choices}

Predictive models for irregular clinical time series differ first in what they consider an input.  One family regularizes observations onto a sequence of time bins and supplies imputed values, masks, or time gaps.  Another works directly with observed events, representing measurement identity, value, and time without filling every unobserved cell.  Within these representations, temporal dependence can be modeled by recurrent state updates, temporal convolution, or attention.  These choices encode distinct assumptions about which temporal relationships should be easy to learn and how the observation process should enter the model \cite{sun_2026_review_irregular_medical_timeseries}.

No representation removes the need to define the task and preprocessing contract.  Dense sequences require a bin width, aggregation rule, and treatment of missing cells.  Sparse events require a variable vocabulary, continuous or discrete encodings of time and value, and a policy for very long event streams.  Positional or temporal encodings provide order information to attention mechanisms, while masks can prevent padded or future positions from entering a computation.  Self-supervised pretraining can exploit unlabeled records, but its utility depends on the pretext task and the relation between pretraining and supervised populations.  The four learned models reviewed below are the predictive families compared by the thesis; their published forms are conceptual precedents rather than proof that every local adaptation is identical.

\subsection{Recurrent baselines: GRU and GRU-D}

The gated recurrent unit (GRU) updates a hidden state sequentially using reset and update gates.  Conceptually, the update gate balances retention of the previous state against incorporation of a candidate state, while the reset gate controls how strongly earlier information contributes to that candidate.  This gating provides a flexible general sequence baseline with fewer gating components than a conventional long short-term memory unit \cite{cho2014gru}.  A standard GRU can learn from a suitably prepared clinical sequence, but it does not inherently know how much time elapsed between observations or why a variable is absent.  Its handling of irregularity therefore depends on the input representation supplied to it.

GRU-D modifies recurrent modeling for multivariate series with missing values.  It uses observation masks and elapsed-time information, and learns decay mechanisms that move missing inputs toward empirical reference values and attenuate hidden states as gaps increase \cite{che2018grud}.  This formulation makes the missingness pattern and time since observation explicit rather than treating imputation as an invisible preprocessing step.  It is one principled response to irregular sampling, not a universal solution: learned decay imposes its own structure, the observation process may differ across sites, and predictive use of missingness does not correct causal bias from measurement decisions.

\subsection{Temporal convolution: TCN}

Temporal convolutional networks apply one-dimensional convolution in temporal order.  Causal convolution prevents future positions from contributing to an earlier output, dilation expands the receptive field without requiring a kernel at every lag, and residual blocks stabilize deeper models.  Because convolutions for many positions can be computed in parallel, a TCN offers a different computational and inductive bias from a recurrent state update \cite{bai2018tcn}.  The canonical TCN study is a general sequence-modeling comparison rather than an ICU-specific method paper.  In an irregular ICU setting, the model still requires an appropriate regular sequence representation, so resampling, aggregation, masks, and time features remain part of its effective design.

\subsection{Sparse-event transformer modeling: STraTS}

STraTS is designed for sparse and irregular multivariate clinical series.  It treats an observed measurement as a triplet of time, variable, and value; continuous embeddings of time and value are combined with a variable embedding, and transformer-style attention contextualizes the observed events.  This avoids constructing a fully imputed variable-by-time matrix and preserves fine-grained observation times.  The method also introduces a self-supervised forecasting objective for pretraining representations before supervised prediction \cite{tipirneni2022strats}.

This design is well matched to sparse event streams, but it does not make preprocessing irrelevant.  Event selection, normalization, truncation, static features, and the variable vocabulary still determine what the transformer sees.  Moreover, the published method and the thesis's supervised multi-label adaptation are distinct objects.  The paper supports the triplet representation, attention mechanism, and pretraining rationale; Chapter~\ref{chap:predictive-modeling-proxy-states} defines the local target heads, training, and export contract.

\subsection{Relationship among the four included model families}

GRU supplies a general recurrent baseline; GRU-D makes missingness and elapsed time explicit within recurrence; TCN provides multiscale temporal filters; and STraTS applies attention to observed event triplets with optional pretraining.  Their suitability can vary with density, preprocessing, labels, sample size, and site-specific observation patterns.  The comparison therefore tests distinct inductive biases rather than a universal sophistication ranking; Chapters~\ref{chap:predictive-modeling-proxy-states} and~\ref{chap:results} define the local implementations and results.

\section{Proxy Phenotyping and Weak Supervision}
\label{sec:background-related-work:proxy-phenotyping-weak-supervision}

\subsection{Electronic phenotyping and rule-derived labels}

Electronic phenotyping identifies patients or records with characteristics of interest from routinely collected EHR data.  Phenotypes may combine diagnosis or procedure codes, laboratory and vital measurements, medications, text, temporal criteria, deterministic rules, statistical models, or machine learning.  The field has consequently developed from manually specified rule sets toward hybrid and learned approaches, while retaining the need to validate the construct for its intended use \cite{banda_2018_electronic_phenotyping}.

Rule-based phenotypes remain attractive when transparency matters: reviewers can trace labels to variables, thresholds, and time conditions and identify implausible clauses.  That inspectability does not make a rule clinically correct.

Transparency is not equivalence to clinical truth.  Thresholds can be sensitive to units and preprocessing, required variables may be missing, and a rule developed at one site may not transport to another.  Temporal ambiguity can cause measurements before, during, and after a clinical episode to be collapsed into one label.  Codes and treatments may reflect billing or care processes as well as physiology.  Rule outputs can therefore be reproducible yet noisy, and a negative label can mean insufficient evidence rather than absence of the condition.  Validation must be tied to the intended use: a cohort-retrieval rule, prediction target, and causal exposure may require different evidence.

\subsection{Clinical-definition sources and approximation}

Clinical consensus documents provide bounded conceptual grounding for organ dysfunction and acute syndromes.  Sepsis-3, KDIGO acute-kidney-injury criteria, the ISTH DIC framework, and the Berlin ARDS definition are multicomponent and context dependent \cite{singer_2016_sepsis3,kdigo_2012_acute_kidney_injury,taylor_et_al_2001_isth_dic,ranieri_et_al_2012_berlin_ards}.

These sources cannot upgrade a partial project rule into a diagnosis.  Local proxies may omit duration, baseline change, imaging, interventions, or context and are only clinically informed by the cited definitions; construct validity requires rule-level review and chart-adjudicated evaluation.

\subsection{Weak supervision and label aggregation}

Weak supervision addresses settings with scarce hand labels and imperfect programmatic sources.  Data programming expresses heuristics as labeling functions and models noisy, conflicting sources to create training data \cite{ratner_et_al_2016_data_programming}.  The relevant lesson is that source quality and correlation matter; adding sources does not guarantee cancellation of shared bias.

A deterministic majority vote is not a learned source-quality model: it applies a fixed rule to binary sources.  Its output is algorithmic aggregation, not clinical consensus, and correlated voters trained on the same rule-derived labels can preserve common error.

\subsection{Implications for this thesis's proxy states}

The thesis uses \emph{clinically inspired proxy state}, \emph{rule-derived proxy label}, \emph{proxy phenotype}, and \emph{weak clinical label} to distinguish analytical constructs from verified diagnoses.  Repository identifiers beginning with \texttt{LAT\_} are historical software names.  They do not establish that a formally identified latent clinical variable exists.  This distinction is essential because the same column interface is used across several stages without making the generating mechanisms equivalent.

Proxy-state evidence spans distinct questions: whether code executes the documented rule, identifiers and columns align, models learn the rule-derived labels, labels correspond to the intended clinical construct, and proxies are temporally coherent causal exposures.  Repository evidence supports implementation, schema, and learnability more directly than clinical construct or causal measurement validity.

The shared proxy interface supplies prediction targets, aggregation inputs, and downstream exposure representations.  It makes handoffs auditable but also links rule, prediction, aggregation, and causal-measurement error; predictive accuracy against rule labels is learnability, not clinical or causal validity.

\subsection{LLM-Assisted Clinical and Causal Knowledge Elicitation}

Advanced LLMs can encode and express substantial medical knowledge under structured evaluation.  In MultiMedQA evaluations, PaLM-family models produced medically relevant answers and reasoning, and instruction tuning improved several measured capabilities \cite{singhal_et_al_2023_llm_clinical_knowledge}.  The same study's human evaluation exposed clinically important gaps, and Med-PaLM remained inferior to clinicians.  This evidence supports the feasibility of using an LLM to generate design proposals; it does not establish reliability for a particular proxy-design task, correctness of every rule, suitability as an autonomous medical expert, or clinical validation of CliniCause outputs.

LLM semantic judgments have also been studied as prior information for causal-graph construction.  Darvariu and colleagues evaluate prompt designs and soft LLM-derived priors for causal discovery, reporting the greatest benefit for edge orientation in selected common-sense benchmark settings while also finding that standalone LLM priors are not uniformly beneficial \cite{darvariu_et_al_2024_llm_causal_graph_priors}.  This is a methodological precedent for treating model output as candidate relations or orientation information, not evidence that the CliniCause DAGs are correct.  The local graphs were not learned from patient-level data, the study did not implement the paper's discovery algorithm, and formal adjustment validity still depends on the correctness and completeness of the project-specified assumptions.

A separate precedent is programmatic labeling: inspectable heuristics generate weak labels for discriminative learning \cite{ratner_et_al_2016_data_programming}.  CliniCause is only conceptually related; its selected decision trees generate deterministic proxy targets, and its later aggregate is a fixed vote rather than a learned noise model.

Within this thesis, the LLM-assisted elicitation protocol is consequently treated as a substantive design-time component with bounded authority.  It produced proposals for proxy-state ontologies, binary decision rules, missingness considerations, and dataset-specific DAGs; the project selected designs and encoded them in deterministic source files.  The prompt artifacts document that process, whereas active tagger and graph scripts define implemented behavior.  Neither the literature above nor the local prompt record establishes clinical validity, complete human review, or causal identification.

\section{Causal Diagrams, Identification, HTE, Overlap, and Sensitivity}
\label{sec:background-related-work:causal-diagrams-identification-hte-overlap-sensitivity}

\subsection{Observational causal questions and target-trial thinking}

Causal analysis begins with a question, not an estimator.  Target-trial thinking makes the desired comparison explicit by specifying eligibility, treatment or exposure strategies, assignment, time zero, follow-up, outcome, causal contrast, and analysis plan before observational data are analyzed \cite{hernan_robins_2016_target_trial}.  This discipline exposes mismatches such as using post-baseline information to define eligibility, comparing prevalent with incident exposure, or allowing eligibility and follow-up to begin at different times.  ICU EHR data make these issues acute because clinical state, treatment, measurement, and selection evolve rapidly and are documented through care processes.

The exposure must also correspond to a sufficiently well-defined intervention or contrast.  Different ways of producing the same measured state can lead to different outcomes, so a contrast on a state such as body mass index---or, here, an illness-state proxy---may violate consistency when the underlying intervention is unspecified \cite{hernan_taubman_2008_well_defined_interventions}.  This concern is not repaired by choosing a flexible estimator.  It requires clarity about what changes, when it changes, and which versions of the exposure are being compared.

Retrospective ICU studies additionally face confounding by indication, changing severity, informative measurement, treatment--confounder feedback, selection, and limitations in outcome measurement.  A review of observational causal inference in intensive care emphasizes careful question formulation, assumptions, data preparation, and reporting across the analysis workflow \cite{smit_2023_causal_inference_icu_scoping_review}.  Marginal structural models provide important background for time-varying treatment and confounder processes, but they are not the point/binary exposure estimator implemented in this thesis.  Nor does this thesis claim a complete target-trial emulation.

\subsection{DAGs, backdoor reasoning, and adjustment}

A directed acyclic graph (DAG) represents assumed causal relations through nodes and directed edges.  Paths into an exposure can encode confounding routes, while colliders require special care because conditioning on them can open otherwise blocked associations.  Backdoor reasoning uses the graph to identify covariates that block noncausal paths without adjusting for inappropriate descendants or colliders \cite{pearl_1995_causal_diagrams}.  A DAG thus makes assumptions inspectable and connects scientific reasoning to adjustment-set selection.

Graphical transparency must not be mistaken for empirical proof.  A DAG does not establish that its arrows are correct, and a set that blocks paths in the drawn graph does not guarantee exchangeability in the clinical population.  Required confounders may be unavailable in the observed data; measurement can be an imperfect proxy for a graph node; and temporal order must be justified scientifically rather than inferred from an arrow.  In this thesis, project-specified DAGs guide observed-variable adjustment at a high level.  Their dataset-specific nodes, mappings, and implementation are deferred to Chapter~\ref{chap:causal-graph-effect-estimation-methodology}.

\subsection{Double machine learning}

Double machine learning (DML) separates the causal parameter of interest from flexible nuisance functions, such as outcome regression and exposure propensity.  Machine-learning estimators can fit these functions, but naively substituting regularized predictions into a causal estimating equation can transfer regularization and overfitting bias to the target parameter.  DML addresses this problem through Neyman-orthogonal scores, which reduce first-order sensitivity to nuisance-estimation error, and cross-fitting, which estimates nuisance functions on data separate from the observations used to evaluate the score \cite{chernozhukov2018dml}.

In a residualization interpretation, covariate-explained components are removed from the outcome and exposure, and the residual association enters the orthogonal score.  Cross-fitting rotates training and evaluation folds so observations can contribute efficiently without using the same fitted nuisance error twice.  These devices broaden the nuisance models that can be used while retaining inference under stated regularity and rate conditions.  They do not eliminate the substantive assumptions of causal inference.  DML cannot correct an omitted confounder that was never measured, create overlap where none exists, turn a proxy into an accurately measured exposure, repair an incorrect DAG, or define an otherwise ambiguous intervention.

EconML provides software implementations for heterogeneous-effect estimators, including DML-related workflows, and is useful for documenting implementation context \cite{oprescu_et_al_2019_econml}.  The software reference is distinct from the theoretical basis: library availability does not validate local nuisance models, cross-fitting behavior, or identification assumptions.

\subsection{Causal forests and heterogeneous effects}

Outcome prediction estimates how outcomes vary with covariates under the observed data distribution; effect estimation asks how an outcome contrast varies between exposure conditions under causal assumptions.  A strong prognostic variable need not be an effect modifier, and a subgroup with high baseline risk need not receive greater benefit from an intervention.  Keeping prognostic and treatment-effect heterogeneity separate is therefore fundamental.

Causal forests adapt random-forest ideas to heterogeneous treatment-effect estimation.  Recursive partitioning seeks regions with different treatment contrasts, while aggregation stabilizes local estimates.  Under unconfoundedness and appropriate forest construction, the method supports conditional-effect estimation and inference \cite{wager2018causalforest}.

These methods are attractive when effect modification is nonlinear or involves interactions, but flexibility does not validate discovered subgroups.  Estimated conditional effects can be unstable in sparse regions, depend on nuisance models and tuning, and identify patterns that do not transport.  Forest feature importance describes the fitted model rather than a biological mechanism.  High-dimensional individual estimates are not automatically clinically actionable.

Critical-care heterogeneity, confounding, informative observation, treatment selection, and covariate shift make conditional effects difficult to validate; a held-out factual outcome cannot reveal an individual's counterfactual error.  ICU causal-analysis guidance therefore emphasizes question formulation, assumptions, data preparation, and reporting rather than estimator flexibility alone \cite{smit_2023_causal_inference_icu_scoping_review}.

CausalPFN is a prior-fitted transformer that amortizes causal-effect estimation through in-context inference after training on simulated data-generating processes satisfying ignorability \cite{balazadeh2025causalpfn}.  It remains exploratory here because the archive does not identify the exact producing package version or checkpoint and provides no DML-equivalent uncertainty, sensitivity, or permutation diagnostics.

\subsection{Overlap and empirical support}

Positivity requires that the exposure conditions being compared are possible within relevant covariate strata.  Its empirical counterpart, overlap or common support, asks whether comparable exposed and unexposed observations exist in the analyzed sample.  When propensity scores approach their boundaries, estimates rely on a small number of observations or extrapolation, producing instability and changing the population for which a comparison is credible.  Work on limited overlap formalizes the variance and target-population consequences and motivates explicit attention to the region of adequate support \cite{crump_et_al_2009_limited_overlap}.

Overlap is especially important for heterogeneous-effect estimation because local estimates divide the sample into increasingly specific covariate neighborhoods.  A global mixture of exposed and unexposed records can coexist with local support failures.  Matching rates, failures, and distances can provide indirect evidence about available comparisons, but they depend on the matching representation and algorithm.  They do not replace propensity-score distributions, common-support plots, or covariate-balance assessment.  Conversely, empirical overlap does not establish exchangeability or validate the DAG.  This thesis therefore treats support diagnostics as necessary qualification and does not claim that positivity has been established.

\subsection{Omitted-variable sensitivity and diagnostic workflows}

Even after adjustment, observational estimates can be biased by variables that affect both exposure and outcome but are absent from the model.  Omitted-variable-bias sensitivity analysis asks how strongly such a variable would need to relate to the residualized exposure and outcome to change a conclusion.  Robustness values provide a compact strength-of-confounding summary, while comparison with observed covariates can make a hypothetical confounder scale more interpretable \cite{cinelli_hazlett_2020_sensitivity}.  Extensions to causal machine learning formulate analogous bias bounds and inference for broader target parameters and flexible nuisance estimation \cite{chernozhukov_et_al_2026_ovb_causal_ml}.

Sensitivity results remain model- and scale-dependent.  Benchmarking is persuasive only when the selected observed covariate is scientifically meaningful, and no observed benchmark necessarily bounds every unmeasured process.  A robustness value does not prove that unmeasured confounding is absent.  Sensitivity analysis also does not validate a DAG, create support, or repair exposure misclassification.  Its role is to quantify how conclusions would respond to specified departures from an identifying assumption.

\subsection{Positioning of the present study}

The literature supplies complementary components rather than the observational evaluation object itself.  Irregular-series models address representation; phenotyping and programmatic labeling create inspectable but imperfect analytical variables; bounded LLM work motivates design proposals; DAGs expose assumptions; DML and forests estimate flexible contrasts; and overlap and sensitivity methods reveal fragility.  CliniCause integrates these components while keeping representation decisions, handoffs, estimator interfaces, diagnostics, and provenance explicit.  Integration makes construction auditable but can also compound rule, prediction, aggregation, graph, support, confounding, and lineage errors.  Chapters~\ref{chap:problem-definition-study-design}--\ref{chap:experimental-design} define the implemented objects and boundaries; none of the cited literature validates the local proxies, graphs, estimates, or causal truth.
